\section{Mobilní komunikace}
	\subsection{Blokové schéma digitálního radiokomunikačního systému, popis bloků, přenosová kapacita systému}
		V blokovém schématu se nachází 3 hlavní části, a to vysílač, přenosový kanál a přijímač.
		\subsubsection{Blokové schéma digitálního radiokomunikačního systému, popis bloků}
			\begin{itemize}
				\item \textbf{Zdroj signálu:} Zdrojem může být například mikrofon, snímač kamery či digitální zdroj jako je počítač.
				\item \textbf{Převod A-D:} Tato číst není nutná, pokud je zdroj signálu digitální. Slouží k digitalizaci analogových signálů.
				\item \textbf{Kodér zdroj:} Slouží k potlačení redundance a irelevance (komprese).
				\item \textbf{Kodér kanálu:} Přizpůsobuje data k přenosu. Dochází k zabezpečení (umělé navýšení redundance) a k prokládání (interleaving).
				\item \textbf{Modulátor:} Signál je namodulován na nosnou. Může se jednat o amplitudovou, frekvenční či fázovou modulaci, případně jejich kombinaci (QAM).
				\item \textbf{Rádiové prostředí:} Fyzikální prostředí, kterým se signál šíří. Nacházejí se v něm náhodné (rušení) a nenáhodné (šum AWGN) jevy.
			\end{itemize}
			\begin{figure}[h]
				\centering
				\includegraphics[width=1\linewidth]{prenosovy_kanal}
			\end{figure}
		\subsubsection{Přenosová kapacita systému}
			Přenosová kapacita systému (Shannon-Hartleyův vztah) je dána rovnicí $C = B \cdot \log_2(1 + \frac{S}{N})$, kde
			\begin{itemize}
				\item C - maximální teoretická rychlost přenosu informace v idealizovaném prostředí,
				\item B - šířka kanálu [Hz],
				\item S - střední hodnota přijímaného výkonu signálu na (přijímač),
				\item N - střední hodnota přijímaného výkonu šumu na (přijímač).
			\end{itemize}
	\subsection{Základní způsoby zpracování signálů v systémech mobilních komunikací (zdrojové kódování, kanálové kódování, prokládání a digitální modulace)}
		\subsubsection{Zdrojové kódování}
			Hlavním úkolem zdrojového kodéru je redukce redundance a irelevance.
			\begin{itemize}
				\item \textbf{Redukce redundance:} Redundantní část je část, která je nadbytečná, je teda možné její množství redukovat. Redukce redundance je vratný proces, nedochází tedy ke ztrátě informace.
				\item \textbf{Redukce irelevance:} Irelevance je část infromace, která není nikterak významná a lze ji vynechat (nepřenášet). Pro výběr, co zachovat a co ne, jsou využívány modely lidského zraku a sluchu, díky kterým lze vyhledat části informace, které člověk nezpracuje či jejich ztráta bude mít minimální vliv na vjem. Jedná se o ztrátový proces.
				\item \textbf{Kompresní poměr (CR - Compress Ration):} Slouží pro kvantitativní posouzení procesu kódovaní a je definován jako poměr vstupu k výstupu: $CR = \frac{R_{vst}}{R_{výst}}$.
			\end{itemize}
		\subsubsection{Kanálové kódování, prokládání}
			Kanálové kódovíní přizpůsobuje přenášený signál přenosovému kanálu. Dochází k zabezpečení přenášených dat a prokládání.
			\begin{itemize}
				\item \textbf{Zabezpečení:} Zabezpečit přenášený data lze pomocí umělého přidání informací (redundanci). Lze využít kódy detekční, které jsou schopné chyby pouze detekovat, nebo korekční, které jsou schopny nejen detekce, ale také opravy určitého množství chyb. Sledovaným parametrem je BER (Bit Error Rate) udávající počet chybně přenesených bitů oproti všem.
				\item \textbf{Prokládání:} Proces přeházení pořadí bitů/bytů (interleaving) na straně vysílače a návratu zpět na správnou pozici (deinterleaving) na straně přijímače. Díky prokládání je shluková chyba (burstová), které vznikla v přenosovém kanálu, rozprostřena na vícero menších chybných úseků, které jsou schopny detekční/korekční zabezpečovací kódy detekovat/opravit.
			\end{itemize}
			Kódy mohou být konvoluční, blokové (Cyklické, Turbo, Reed-Solomonovi, BCH, LDPC). Lze je také zřetězit za sebe pro lepší zabezpečení. Mezi jednotlivé kodéry jsou vloženy prokladače.
		\subsubsection{Digitální modulace}
			Modulace je proces, při kterém dochází k ovlivnění alespoň jednoho z parametrů nosné vlny v závislosti na modulačním signálu, a to amplitudy, frekvence či fáze.
			\begin{itemize}
				\item \textbf{Amplitudová:} Při této modulaci je modulačním signálem ovlivňována amplituda nosné vlny. Maximální změnu (snížení) amplitudy nosné modulovaným signálem označujeme jako hloubku modulace.
				\item \textbf{Frekvenční:} Vysílaná frekvence se mění v závislosti na modulačním signálu. Maximální rozdíl frekvence nosné a modulované nosné $\delta f$ nazýváme frekvenční zdvih.
				\item \textbf{Fázová:} Dochází k ovlivnění (změně) fáze nosné. Příklad je doustavová modulace DPSK, u které je použita fáze nosné 0\textdegree a 180\textdegree.
			\end{itemize}
			Mezi základní parametry diitálních modulací patří
			\begin{itemize}
				\item Přenosová rychlost (bitová)
				\item Symbolová rychlost
				\item Bitová chybovost (BER)
				\item Symbolová rychlost (SER)
				\item Vektorová chyba (EVM) - menší hodnota - menší úroveň šumu v konstelačním diagramu [\%]
				\item Energetická účinnost
				\item Spektrální účinnost [bit/s/Hz]
			\end{itemize}
			Zástupci:
			\begin{itemize}
				\item QPSK
				\item O-QPSK (Offset) - neprochází 0
				\item $\pi$/4–DQPSK - diferenční modulace (kombinace lde změny nsoné, nikoli cílového bodu v konstelačním diagramu)
				\item MSK (Minimum Shift Keying) - realizace - VCO
				\item GMSK (Gaussian-filtered MSK) - potlačení postranních laloků
				\item OFDM (Orthogonal Frequency Division Multiplex) - vícero nosných uspořádaných tak, aby nedocházelo ke vzájemnému ovlivňování
			\end{itemize}
	\subsection{Systémy s mnohonásobným přístupem a metody multiplexování, proces rozprostírání signálu, výhody a nevýhody jednotlivých přístupů}
		\subsubsection{Systémy s mnohonásobným přístupem a metody multiplexování}
			Jedná se o přístupy, díky kterým je možné sdílet rádiové prostředí vícero účastníky.
			\begin{itemize}
				\item \textbf{FDMA:} Jedná se o frekvenční dělení. Každý z účastníků ma vyhrazenou část spektra pro komunikaci.
				\item \textbf{TDMA:} Každému účastníkovi je přidělen časový interval (time slot), ve kterém může realizovat přenos.
				\item \textbf{CDMA:} Komunikace vícero účastníků probíhá ve stejném kanále, každý však využívá svůj jedinečný pseudonáhodný kód. Využívá buďto přímou modulaci DS (Direct Sequence) nebo kmitočtové skákání FH (Frequency Hopping), případně Hybridní DS-FH.
				\item \textbf{ALOHA:} náhodný přístup
			\end{itemize}
		\subsubsection{Proces rozprostírání signálu}
			Úzkopásmový signál lze rozprostřít na širší nízkoúrovňový signál SS (Spread Spectrum). Pro rozprostírání je generována pseudonáhodná posloupnost PNP čipovou rychlostí (CR - Chip Rate) výrazně vyšší, než je bitová rychlost modulačního signálu. PNP má ideálně impulzní autokorelační funkci a vzájemně mezi sebou mají PNP nulovou korelační funkci. Rozprostírání probíhá dvěma způsoby:
			\begin{itemize}
				\item \textbf{DS (Direct Sequence):} Nosná je přímo modulována PNP.
				\item \textbf{FH (Frequency Hopping):} PNP ovlivňuje frekvenci nosné. Dle PNP se kmitočet neustále mění.
			\end{itemize}
			\begin{itemize}
				\item \textbf{Výhody:} Hlavní výhodou je odolnost proti rušení (úzkopásmovému i širokopásmovému). Díky nutnosti znalosti PNP pro dekédování dochází také k utajení přenášené informace.
				\item \textbf{Nevýhody:} Realizace funkčně i obvodově náročná (synchronizace a časování).
			\end{itemize}
		Nejčastěji se využívají kombinace přístupů (FDMA-TDMA (GSM), CDMA-TDMA (UMTS)).
	
	\subsection{Způsoby přenosu (simplex, duplex) a typy spojování (komutované a paketové), diverzitní příjem, diverzitní systémy}
		\subsubsection{Způsoby přenosu (simplex, duplex) a typy spojování (komutované a paketové)}
			Přenos je možné realizovat vícero způsoby.
			\begin{itemize}
				\item Způsob přenosu
				\begin{itemize}
					\item \textbf{Simplex:} Jde o jednosměrný přenos bez zpětného kanálu. Zástupcem je například TV vysílání či rádio.
					\item \textbf{Half-duplex:} Poloduplex spočívá ve střídání směrů komunikace. Střídání však není striktně vynucováno. Konec vysílání bývá často doprovázeno předem dohodnutým znakem. Příkladem jsou vysílačky (\uv{přepínám}).
					\item \textbf{Duplex:} Obousměrný přenos realizovaný ve stejný okamžik. Může se jednat o duplex kmitočtový (FDD) využívající rozdílné frekvence pro každý směr či časový duplex (TDD) využívající časové úseky (time sloty) pro střídání směru komunikace.
				\end{itemize}
				\item Způsob spojování
				\begin{itemize}
					\item \textbf{Komutovaný:} Spojení s přepojováním okruhů. Před samotnou komunikací je nutné sestavit komunikační okruh, který je během celé komunikace rezervován, i když je nevyužíván.
					\item \textbf{Paketový:} Tok dat je rozdělen na menší části (pakety). Každý paket je opatřem záhlaví se zdrojem, cílem, prioritou, ... a je samostatně směrován sítí. Paketové přenosy jsou realizovány buďto službou bez spojení, u které není garantováno ani potvrzováno doručení, nebo službou se spojením, při které je vytvořen virtuální okruh.
				\end{itemize}
			\end{itemize}
		\subsubsection{Diverzitní příjem}
			Úroveň signálu, vlivem okolí, překážek nebo pohybu vysílače či přijímače, může kolísat. Takové kolísání je označováno jako únik (fading). Řešením může být právě diverzitní příjem. Jedná se o techniku příjmu omezující (potlačující) rušivé vlivy. Předpokladem je minimální vzájemná korelace úniků v jednotlivých kanálech. Při splnění techto podmínek je možné přijímat přímé i odražené signály. Dělí se na diverzitní systém s výběrem:
			\begin{itemize}
				\item \textbf{Prostorový:} Na přijímací straně je daný počet samostatných antén s rozestupem několika vlnových délek $\lambda$, minimálním však $\lambda/2$.
				\item \textbf{Úhlový:} Využívají se směrové antény do různých směrů (příjem odrazů).
				\item \textbf{Polarizační:} Dvě antény s rozdílnými polarizacemi (horizontální x vertikální).
				\item \textbf{Kmitočtový:} Systém, který může využívat jednu anténu, avšak počet vysílačů a přijímačů je dán počtem použitých kmitočtů. Frekvenční odstup by měl být 2-4\% kmitočtu nosné.
				\item \textbf{Časovým:} Systém využívající rámce TDM.
			\end{itemize}
			Dále se dají dělit dle počtu antén na:
			\begin{itemize}
				\item \textbf{SISO (Single Input Single Output):} Systém s jednou anténou na vysílání i příjem.
				\item \textbf{SIMO (Single Input Multiple Output):} Jedna vysílací a vícero přijímacích antén.
				\item \textbf{MISO (Multiple Input Single Output):} Vícero vysílacích antén a jedna přijímací anténa.
				\item \textbf{MIMO (Multiple Input Multiple Output):} Vícero vysílacích i přijímacích antén.
			\end{itemize}
	\subsection{Struktura obecné buňkové sítě a její hlavní výhody, procedura "handover" a její různé typy}
		\subsection{Struktura obecné buňkové sítě a její hlavní výhody}
			Buňková struktura je postavena na myšlence rozdělení daného území na menší oblasti (buňky). V každé buňce je umístěna základnová staníce (BTS). Díky vzdálenosti a nižším vysálacím výkonům je možné opakovat kmitočty již ob buňku, je tedy teoreticky možné s využítím omezeného kmitočtového pásma pokrýt nekonečně velké území. Pro zvýšení kapacity je možné buňky dále dělit, provést sektorizaci a také provádět dynamické přidělování kanálů (dle okolních buňek). Celý proces návrhu a měření je označován jako kmitočtové plánování.\\
			Buňky lze dělit dle velikosti na:
			\begin{itemize}
				\item \textbf{Femtobuňky:} Poloměr je v desítkách metrů. Vykrývají malá nepokrytá místa (domácnosti, kanceláře). Jsou označovány také jako \uv{deštníkové (umbrella cells)}.
				\item \textbf{Pikobuňky:} Poloměr do 50 m. Určena pro místa s vysokou koncentrací lidí (nádraží, obchodní domy).
				\item \textbf{Mikrobuňky:} Poloměr do 1 km. Oblasti s vyšším provozem (centra měst).
				\item \textbf{Makrobuňky:} Poloměr až desítky km. Oblasti s malou hustotou provozu (venkovské oblasti).
			\end{itemize}
		\subsection{Procedura "handover" a její různé typy}
			Při pohybu či zhoršení podmínek v daném přenesovém kanále nastane situace, ve které je nutné přepojení na jinou základnovou stanici. Tento proces se nazývá \uv{handover}.\\
			Dle způsobu přepojení se může jednat o handover:
			\begin{itemize}
				\item \textbf{Tvrdý (hard):} Při tomto typu dojde k odpojení od první BTS a poté teprve k připojení ke druhé. Dochází ke krátkému (asi 10 ms) výpadku spojení.
				\item \textbf{Bezešvý (seamless):} Před odpojení od původní BTS se mobilní stanice připojí k nové BTS. Až po úspěšném připojení k nové BTS dojde k odpojení od staré BTS. Na krátký okamžik je tedy připojena k oběma BTS a nedochází k výpadku.
				\item \textbf{Měkký (soft):} Mobilní stanice je připojena k vícero BTS zaráz. Díky těmto paralelním spojení je možné přepojování mezi různými BTS bez výpadku, jelikož je vždy alespoň jedno spojení aktivní.
			\end{itemize}
			Dle toho, která část měří kvalitu spojení, může být:
			\begin{itemize}
				\item \textbf{Sítí řízený handover (NCHO – Network Controlled Handover):} Měření kvality provádí BTS. Na základě měření provádí rozhodnutí o přepnutí. Na mobilní stanice nejsou kladeny žádné speciální požadavky.
				\item \textbf{Handover řízený mobilní stanicí (MCHO – Mobile Controlled Handover):} Měření kvality všech kanálů provádí mobilní stanice i BTS. O přepnutí rozhoduje mobilní stanice. Rozhodnutí předá systému a ten zajistí samotné přepojení.
				\item \textbf{Sítí řízený handover s asistencí mobilní stanice (MAHO – Mobile Assisted Handover):} Mobilní stanice monitoruje signály okolních BTS a výsledky předává BTS, ke které je aktuálně připojena. Současně mobilní stanice i BTS provádí měření aktuálního vzájemného spojení. Na základě naměřených údajů systém provádí rozhodnutí o přepojení i jeho realizaci.
			\end{itemize}